%!TEX root = ../../csuthesis_main.tex
\chapter{引言}

持续学习旨在使模型能够在数据分布随时间发生变化的情况下不断接受新信息的同时尽可能地保持之前学到的知识~\citep{wang2024survey,zhou2024class,zhou2023cil_survey_zh,wang2025llm_cl_zh,leo2024survey}。这与传统的静态监督学习有很大的区别，在传统的静态监督学习中可以一次性获取所有的训练样本，而在持续学习中是要面对一个一个到来的数据流、不断增加的类别数目以及资源有限的情况等。因此问题也就随之改变，不再是单纯的“如何拟合当前的数据”，而是要解决“如何在记住以前的任务的情况下还要学习新的任务”。
~\citep{wang2024survey,mccloskey1989catastrophic,goodfellow2013catastrophic}。这个问题在医学图像分析~\citep{yang2023medmnist,derakhshani2022lifelonger}、语义分割~\citep{cermelli2020mib,douillard2021plop}、时间序列感知~\citep{qiao2024tscil,huo2025continue}以及多模态大语言模型~\citep{zhai2023investigating,wu2024continual,shi2024continual,guo2025mcitlib}等多个方面都尤为突出——即一方面需要模型具有长期学习的能力，同时在鲁棒性、有效性和隐私性等方面也有着较高的期望。

\section{研究背景与意义}

从 ResNet~\citep{he2016resnet} 到 ViT~\citep{dosovitskiy2020image}、再到当前的多模态大模型~\citep{liu2023llava,chen2024internvl}，深度学习在图像分类、目标识别~\citep{radford2021learning}、语义分割~\citep{chen2017deeplabv3,wang2019dgcnn}、序列建模~\citep{qiao2024tscil}等方向上的进展，构成了过去十几年这一领域最显眼的部分。这里有一个常常被默认、但实际上很强的前提：模型是在训练阶段可以一次性看到全部数据"的条件下被训练出来的~\citep{wang2024survey,zhou2024class}。

但是现实世界并非如此，实际上，在临床诊断中，新型疾病、新病理亚型、新成像设备等不断加入到数据库中~\citep{yang2023medmnist,derakhshani2022lifelonger,sun2023few}；机器人或者自动驾驶车辆上线后，会遇到之前未见过的物体类型~\citep{rebuffi2017icarl,hou2019learning,zhou2023cil_survey_zh}；传感器端时间序列任务，会不断出现新的人类活动或者新的工作情况~\citep{qiao2024tscil,huo2025continue}；而对于多模态大语言模型来说，这更是司空见惯的情形——新的任务进入系统几乎不可避免，而且需要相应的引入新的技能模块~\citep{zhai2023investigating,wu2024continual,shi2024continual,zhao2025mllm,guo2025mcitlib}。综合上述这些情况可以看出，无论哪种情况都指向一个共同点，即模型要在无法充分回顾过往信息的情况下，接受一个又一个的新任务。

持续学习中最棘手的问题是灾难性遗忘——它的具体表现是模型在学习新任务的过程中，参数更新破坏了旧任务上已经组成的判别边界，进而激发早期任务性能的明显下滑~\citep{goodfellow2013catastrophic,mccloskey1989catastrophic,kirkpatrick2017ewc,zenke2017continual}。如果暂时把数据可访问性这一约束放在一边，把各阶段数据 $\mathcal{D}_1,\ldots,\mathcal{D}_t$ 同时呈现给模型，那么理想情况下的最优参数可由如下联合训练目标给出：

\begin{equation}
\theta^{\star}_{1:t} \;=\; \arg\min_{\theta}\sum_{k=1}^{t}\mathbb{E}_{(x,y)\sim\mathcal{D}_k}\big[\ell\big(f_\theta(x),y\big)\big],
\label{eq:cl_oracle}
\end{equation}
该目标在持续学习领域被视为上界参考~\citep{wang2024survey,zhou2024class}。但在真实的持续场景中，阶段 $t$ 只能观察到 $\mathcal{D}_t$，所以反向传播带来的更新只映射当前任务的梯度方向 $g_t$。一旦某个历史任务的代表梯度 $g_k$（$k<t$）与之夹角超过九十度，即符合
\begin{equation}
g_t^{\top}g_k\;<\;0,
\label{eq:grad_conflict}
\end{equation}
当前的更新就会同时把旧任务损失变高，而这正是遗忘在几何层面上的成因~\citep{goodfellow2013catastrophic,lopez2017gradient,chaudhry2019continual,zhai2023investigating}。图~\ref{fig:catastrophic_forgetting} 对此给出了非常直观的刻画：损失景观表明参数自历史任务最优点 $\theta^{\star}_1$ 沿当前任务的梯度方向漂移到 $\theta_t$ 之后已经偏离了原来的低谷，所以由此引致旧任务性能下降，导致了遗忘。

\begin{figure}[ht]
\centering
\includegraphics[width=0.85\textwidth,height=0.4\textheight,keepaspectratio]{images/灾难性遗忘与梯度冲突示意图.png}
\caption{灾难性遗忘与梯度冲突示意图。左：在损失景观中参数沿新任务梯度方向移动，致使在旧任务最优点 $\theta^{\star}_1$ 附近的性能下滑；右：当前任务的梯度 $g_t$ 与历史任务的梯度 $g_k$ 夹角超过九十度，满足内积条件 $g_t^{\top} g_k < 0$，从几何层面揭示了遗忘的来源。}
\label{fig:catastrophic_forgetting}
\end{figure}

对于这些问题，根据不同的更新机制，已经有一些相关的工作可以分为三类传统的做法即基于回放、基于正则化以及基于结构扩展的方法~\citep{wang2024survey,zhou2024class,zhou2023cil_survey_zh,wang2025llm_cl_zh}，第一类是回放方法，它们的基本思想就是利用保存的历史样本、缓存特征或者伪样本来减少遗忘，例如iCaRL~\citep{rebuffi2017icarl}、EEIL~\citep{castro2018end}、ER~\citep{rolnick2019experience}、DER~\citep{buzzega2020dark}和GEM/A-GEM~\citep{lopez2017gradient,chaudhry2019continual}等，这些方法在公开的数据集上效果不错，但是需要大量存储历史的信息，在医疗领域、隐私敏感的信息或者是嵌入式设备中这样的假设一般是不能满足的；第二类是正则化方法，它们希望通过限制重要的参数变化来保护之前的知识，如EWC~\citep{kirkpatrick2017ewc}、SI~\citep{zenke2017continual}、MAS~\citep{aljundi2018memory}和LwF~\citep{li2017learning}等，但是这种做法也有缺点，那就是参数的重要程度一般只能通过局部近似得到，而且随着任务的数量增加，正则化的惩罚会越来越难平衡稳定性和灵活性~\citep{zhou2023cil_survey_zh}；第三类是结构扩展或者参数隔离方法~\citep{rusu2016progressive,aljundi2017expert,wang2022dualprompt}，它们试图通过引入可动态调整网络、专家模块或是任务相关的参数降低不同任务间的干扰，但同时也面临着如何设计路径、选取何种组件等问题~\citep{hu2021lora,wang2023orthogonal,chen2024coin}。

进一步地，大部分持续学习方法都基于梯度下降~\citep{wang2024survey,zhou2024class}。新任务产生的梯度直接作用到已有的参数上，而且任务之间的梯度方向经常是不同的，这也是造成遗忘的一个原因~\citep{goodfellow2013catastrophic,lopez2017gradient,zhai2023investigating}。也就是说，在使用反向传播进行更新的情况下，旧的知识就会受到新的知识的影响。因此，从优化的角度来研究持续学习问题是有意义的。与基于梯度的方法不同，闭式解方法以解析的形式求出最优分类器参数而无需进行多次反向传播和数据回放\citep{zhuang2022acil,zhuang2024dsal,fang2024air,GKEAL_Zhuang_CVPR2023,GACL_Zhuang_NeurIPS2024}。因此它们是很有潜力实现隐私保护并且高效持续学习的方法~\citep{zhuang2022acil,zhuang2024dsal,ZHUANG2024107285,he2024realrepresentationenhancedanalytic}。

闭式解持续学习的研究意义可以从多个方面进行阐述，在理论上，它提供了一个持续学习能否不需要梯度更新的新视角去研究~\citep{zhuang2022acil,zhuang2024dsal,fang2024air}，同时也是一个从理论上研究稳定性和可塑性一致性的切入点；在算法上，因为闭式解一般可以把持续更新转化为一个矩阵递推的过程，所以只需要一轮或者一次的数据前向传播就可以实现对新任务的学习从而大大降低了计算成本~\citep{GKEAL_Zhuang_CVPR2023,GACL_Zhuang_NeurIPS2024,AEF-OCL_Zhuang_TVT2024}；从应用的角度来看，该方法本身就很适合那些不能长期保存原始样本的应用场景如医学图像、用户个人隐私信息或工业物联网终端设备等；而在系统层面，由于闭式解具有较好的灵活性它可以被应用于诸如疾病的识别和时间序列分类这类较为传统的分类任务~\citep{rebuffi2017icarl, qiao2024tscil}，也可以推广到语义分割以及多模态大语言模型专家路由这样更为复杂的问题中去~\citep{cermelli2020mib, douillard2021plop, zhao2025mllm, chen2024coin}。

\section{国内外研究现状}

关于持续学习相关工作，目前已有的研究主要分为三个流派，即回放、正则化、结构扩展~\citep{wang2024survey,zhou2024class,zhou2023cil_survey_zh,wang2025llm_cl_zh,leo2024survey}。最早被广泛关注的是分类任务~\citep{rebuffi2017icarl,kirkpatrick2017ewc,li2017learning}；随后该领域逐渐拓展到分割以及时序建模等方面~\citep{cermelli2020mib,douillard2021plop,qiao2024tscil,huo2025continue}；近年来又进一步推广至多模态大语言模型方面~\citep{zhai2023investigating,wu2024continual,shi2024continual,guo2025mcitlib,zhao2025mllm}。



近年来，解析学习以及闭式解的方法逐渐受到重视~\citep{zhuang2022acil,zhuang2024dsal,fang2024air,GKEAL_Zhuang_CVPR2023,GACL_Zhuang_NeurIPS2024,ZHUANG2024107285,he2024realrepresentationenhancedanalytic}。它们的基本出发点并不是通过多次迭代来改进分类头，而是将特征映射后的学习问题转化为最小二乘或者岭回归的形式，然后利用矩阵求逆或者递归来获得一个最优解。这种方法最初主要是用于单次批学习中，在此基础上发展出了适应不断有新样本到来情况下的递推求解方法，在无需访问旧的原始样本的情况下也能保存下来之前的信息并且进行准确更新~\citep{zhuang2022acil,zhuang2024dsal,fang2024air,AEF-OCL_Zhuang_TVT2024}。本文工作也是沿着这个方向进行研究，“历史知识保留”不再是通过重播或者梯度限制的方式实现，而是转变成对于自相关矩阵以及互相关矩阵的递推维护的问题，从而使得整个过程更加简洁明了。

然而，闭式解持续学习还在快速发展之中，目前的研究大多局限于图像分类这类比较常见的任务设置~\citep{zhuang2022acil,GKEAL_Zhuang_CVPR2023,fang2024air}。而对于更加接近实际应用场景的问题，还存在许多亟需解决的问题，如如何将闭式解的思想从分类推广到像素级或者点级别的预测，比如语义分割~\citep{cermelli2020mib,douillard2021plop}；如何在类内变化较大的数据集例如时间序列上实现稳定性和可塑性并存的问题~\citep{qiao2024tscil,huo2025continue}；如何将闭式解的方法与参数分离或者专家系统结合起来应用于多模态的大语言模型的路由问题~\citep{hu2021lora,wang2023orthogonal,chen2024coin,zhao2025mllm}；如何设计一种通用的数学结构使得各种不同类型的任务都可转化为一类递归解析问题进行处理。对这些问题的研究不仅可以促进闭式解持续学习的发展和完善，也可以使它更好的服务于各种不同的领域。

\section{本文研究内容}

本文围绕"基于闭式解算法的持续学习方法与算法完成"这一主题，在不同任务场景上展开成体系的研究。整体的研究思路如下：首先在数学层面构建统一的递归岭回归闭式解框架，将其作为全文方法论的核心；其次以该框架作为基础，分别面向医学图像疾病分类、时间序列分类、语义分割以及多模态大语言模型专家路由四类典型任务给出具体的应用方法；最后在统一的实验体系下对上述方法做整体验证。

本文的具体研究内容如下。

第一，本文构建了递归岭回归闭式解算法的统一数学基础。凭借冻结编码器、仅对解析分类头做持续更新，持续学习中的训练问题被转化为带正则项的最小二乘问题；在此之上，引入递归更新机制，仅借助上一阶段的统计量与当前任务的数据就能得到与联合训练一致的闭式解。这一部分构成全文方法论的核心。

第二，在医学领域中由于新出现疾病种类不断增加以及以往的数据不能保存的隐私问题，本文以疾病的分类作为应用出发点，将上述统一闭式解推广到医学影像中的分类持续学习的方法并从平均准确度、遗忘率、训练开销及数据隐私性几个方面对其优劣进行了阐述~\citep{yang2023medmnist,derakhshani2022lifelonger,sun2023few}。

第三，对于时间序列上的持续学习问题，我们研究闭式解方法对于类内变化的鲁棒性，给出基于多尺度特征融合以及随机高维映射的时间序列闭式解方案~\citep{qiao2024tscil,qiao2023dt2w,huo2025continue}，并通过集成提高其鲁棒性。

第四，对于语义分割的持续学习任务，我们使用闭式解分类头以及高维随机映射的方法，并对二维图像以及三维点云设计不同的伪标签生成方式~\citep{cermelli2020mib,douillard2021plop,chen2017deeplabv3,wang2019dgcnn}从而得到一种无需样本回放的方式进行类别持续的语义分割方法。

第五，在面向多模态大语言模型的持续学习场景下，我们在专家路由问题上也应用了闭式解的思想~\citep{zhao2025mllm, hu2021lora, wang2023orthogonal, chen2024coin, liu2023llava, chen2024internvl}：采用多模态大语言模型最后一层隐藏状态表示语义并用解析方式进行专家选择头学习从而实现快速并且无遗忘的专家路由。

第六，本文在统一实验体系下对上述方法做了一轮覆盖性的评估，覆盖医学图像、时间序列、二维图像与三维点云语义分割以及多模态大语言模型专家路由这四类任务，从平均准确率、遗忘率与运行效率等多个角度全面证实所提方法的有效性。

本文所做的研究并不是针对一个任务单独进行，而是以闭式解持续学习这一主线贯穿始终，由浅入深，从理论一直到视觉、时序以及多模态的任务，在一个相同的实验设置下进行分析。

\section{本文的主要贡献}

围绕以上六个方面的工作，在此基础上总结出本文的主要创新点共有四个方面：第一个方面是全文方法论思想主线，贯穿这六方面的研究工作即用闭式解代替持续阶段反向传播的整体思路；第二个方面是与第一个方面的第一个研究内容相对应的角度，给出了一个统一的递归岭回归闭式解的框架；第三个方面是与第二个到第五个方面的研究内容相对应的角度，把这个框架推广到了医学图像、时间序列、语义分割以及多模态大语言模型专家路由这四个代表性应用领域；第四个方面是与第六个方面的工作相对应的角度，从效率、隐私和工程可行性三个方面体现出了闭式解的优势。即：

\textbf{思路提炼。}本文总结和归纳了梯度更新是导致持续学习中出现遗忘的主要原因~\citep{goodfellow2013catastrophic,lopez2017gradient,chaudhry2019continual}，然后给出了一种用闭式解形式的方法来代替持续学习中的反向传播的思想，从而给持续学习提供一种不同于回放、正则化以及结构扩展这三种已有的方法的新思路。

\textbf{统一框架。}本文提供一种统一的递归岭回归闭式解框架~\citep{zhuang2022acil,zhuang2024dsal,fang2024air}，利用保存的历史统计量进行迭代更新解析分类头，并从理论上证明这种更新方法得到的结果等价于相应的联合训练结果，有利于后面各种应用场景。

\textbf{多任务推广。}本文将闭式解持续学习方法扩展到多种任务上，包括医学影像疾病的诊断~\citep{yang2023medmnist}、类持续语义分割~\citep{cermelli2020mib,douillard2021plop}、时间序列分类~\citep{qiao2024tscil}以及多模态大语言模型专家路由~\citep{zhao2025mllm,chen2024coin}等，在不同方面验证了这种方法的有效性和鲁棒性。

\textbf{效率及隐私优势。}本文介绍了闭式解方法的优势。相比于需多次训练并进行历史数据回放的方法，本文提出的方法只需要一次遍历当前任务的数据即可完成更新工作，因此更加适合资源有限或者对数据有较高保护要求的应用场景~\citep{de2021continual,GKEAL_Zhuang_CVPR2023}。

\section{论文组织结构}

全文一共分为六章，各章节具体安排如下。

第一章是绪论部分，主要包括本课题研究目的及意义、国内外相关研究概况、本文的研究内容、创新点等以及全文的整体框架等。

第二章对持续学习相关研究进行总结，在问题定义、应用场景及方法分类的基础上概述现有工作，着重介绍基于回放、基于正则化、基于结构扩展或参数隔离以及基于解析学习这四种方法的思想、不足之处及其评价标准等，为后文具体方法做铺垫。

第三章为方法章节，在此部分主要完成递归岭回归闭式解算法的基础理论研究工作即给出问题定义、岭回归闭式解、持续递推公式及其相关证明，这是全文方法论基础部分。

第四章是应用章节，在第三章建立统一数学基础上，对医学图像疾病分类、时间序列分类、二维图像和三维点云语义分割以及多模态大语言模型专家路由四种常见问题提出相应的闭式解持续学习解决方案。

第五章为实验结果及分析，在同一套实验环境下对第四章提出的四种不同类型的任务进行测试。本章首先介绍各个任务所用到的数据集和相关配置信息、所使用的评估标准及比较方式，然后给出主要的结果并对其原因进行讨论。

第六章为结论，总结了本论文主要研究内容并指出现有闭式解持续学习方法在未来可以在更大规模模型上、更为复杂的任务结构上以及开放世界学习中进一步应用和发展。

为了使文章层次分明，在第1至第6章后都有一节小结，对本章内容进行简单总结并且过渡到下一部分内容。

\section{本章小结}

本章首先介绍了持续学习研究的意义及其面临的挑战，在梯度反向传播造成灾难性遗忘的同时还存在隐私和存储成本问题制约了其应用；其次总结了回放、正则化、结构扩展及解析持续学习四种常见方法并提出在闭式解角度上仍有一定的改进余地；最后概述本文工作内容、四大亮点以及全文框架。接下来的第2章将全面介绍各种持续学习的方法并给出本文使用的评估标准。
